\scanpage{4}\mspage{3}
À cause de cette interprétation, le passage de $\widetilde F\subset\mathcal{M}$ à $\operatorname{Fig}_{\mathcal{M}}(F)$ est à tel point évident, que \inserted{la considération de} cette figure ensembliste dans $\mathcal{M}$ me semble \unsure{rarement|vraiment} importante -- mais à voir.....

\deleted{Quand on \unsure{identi}}

On pourra utiliser les deux applications (1.3) pour identifier une figure, suivant les cas, \unsure{à} une partie de $\mathcal{M}$ [via $F\unsure{\mapstochar\relbar|\mapsto}\widetilde F$] ou à une figure ensembliste dans $\mathcal{M}$ [via $F\mapsto\operatorname{\text{Dépl}}(F)$]. Dans le premier cas, \unsure{comparant|composant} avec $\mathcal{M}\unsure{\lhook\joinrel\relbar|\hookrightarrow}\mathcal{F}$ (1.2)\deleted{\unsure{(}}, par la formule \unsure{1,5|1.5}. Cette dernière inclusion s'identifie à
\[
X\mathrel{\mapstochar\relbar\mapstochar}\widetilde X=\{Y\in\mathcal{M}\mid Y\leq X\}
\qquad \mathcal{M}\to\unsure{\mathcal{P}|\mathscr{P}}_{\mathrm{ferm}}(\mathcal{M}).
\]

Donc il faudra être prudent \deleted{\unsure{alors}} \inserted{\unsure{avec}} l'identification dont il a été question, des\deleted{\unsure{une|un}} $X\in\mathcal{M}$ \unsure{ou} multistrates à des figures particulières, quand celles-ci sont \unsure{interprétées} comme des parties de $\mathcal{M}$: La figure associée à $X\in\mathcal{M}$ n'est nullement $\{X\}$, mais $\widetilde X\unsure{\rbrace|,|;}$ \deleted{\unsure{$F$|l'}} \unsure{ens} des multistrates majorées par $\unsure{x|X}$ pour $\leq$ (ou encore, des sous-multistrates, ou multistrates incidentes à $X$, ou encore des \deleted{\unsure{$F$}} \uline{strates} de $x$).
